Urban recovery is usually analyzed through observed degradation and rebound. The central question is whether the system returned to normal, how quickly it did so, and which places experienced the largest or longest deficits. This paper asks a different question: what portion of the observed loss was actually recoverable through intelligent management under realistic constraints?

This framing separates observed resilience from recoverable resilience. A city can recover quickly because endogenous dynamics are strong, even when management decisions have little additional leverage. Another city can appear fragile in observed data while still containing highly recoverable losses if targeted intervention can reduce exposure or restore access.
